% Drafted introduction and related work sections for the ELD manuscript.
% The writing agent must rewrite this prose into a unified voice; do not \input this file.

\section{Introduction}

Spatial landmarks underpin a growing share of quantitative tissue biology. They allow cross-sample comparison of histological heterogeneity, they make it possible to track regions of interest over successive imaging experiments, and they are essential whenever tissue sections must be registered to a common coordinate framework so that molecular measurements from different slices or different modalities can be compared in the same anatomical frame~\cite{moses2022museum,marx2021method}. Spatially resolved transcriptomics platforms such as the barcoded-array Visium assay introduced by St\aa hl and colleagues~\cite{stahl2016visualization} and probe-based in situ sequencing approaches~\cite{ke2013insitu} have dramatically expanded the volume of spatial data that must be aligned, and projects that combine spatial transcriptomics with mass spectrometry imaging~\cite{vicari2024spatial} or that build organ-scale atlases~\cite{asp2019spatiotemporal} push this alignment requirement into a setting where accurate, reproducible landmarks are in short supply. Manual annotation currently fills that gap, but it is labour intensive and quickly becomes the dominant cost of analysis as the number of sections grows.

Deep learning has made automatic landmark detection practical in general computer vision, and there is a natural temptation to transplant these ideas wholesale into spatial omics. Supervised keypoint detectors such as the stacked hourglass network~\cite{newell2016stacked} deliver state-of-the-art accuracy on well-curated benchmarks such as CelebA~\cite{liu2015celeba}, MAFL~\cite{zhang2016learning}, and AFLW~\cite{kostinger2011aflw}, but they depend on large amounts of labelled data that simply do not exist for most histology, transcriptomic, or metabolomic tissue sections. Unsupervised approaches built around landmark equivariance or conditional image generation~\cite{thewlis2017unsupervised,jakab2018conditional,zhang2018unsupervised,sanchez2019object,jakab2020selfsupervised} remove the labelling bottleneck but leave three unresolved issues for tissue data: they are typically trained on tens of thousands of images rather than the five to twenty sections of a spatial omics experiment, they focus on affine deformations whereas real tissue sections warp elastically during sectioning and imaging, and they provide no mechanism for handling simultaneous measurements in different modalities such as H\&E, Visium~\cite{stahl2016visualization}, or mass spectrometry imaging~\cite{vicari2024spatial}.

A second line of work approaches the problem from the registration side rather than the landmark-detection side. Tools such as CODA~\cite{kiemen2022coda} assemble three-dimensional tissue reconstructions from serial sections using tailored registration pipelines, the benchmark framework of Kartasalo and colleagues~\cite{kartasalo2018comparative} quantifies reconstruction accuracy across a panel of algorithms on a mouse prostate dataset, and STalign~\cite{clifton2023stalign} aligns spatial transcriptomics data via Large Deformation Diffeomorphic Metric Mapping. Landmark-based tools such as Eggplant~\cite{bergenstrahle2021eggplant} instead rely on a small number of manually annotated landmarks to project Visium sections onto a common coordinate framework. These methods produce high-quality alignments but either sidestep landmark learning entirely or require an operator to provide the landmarks by hand.

We propose Effortless Landmark Detection (ELD) to bridge these gaps. ELD builds on the unsupervised landmark detector of Sanchez and Tzimiropoulos~\cite{sanchez2019object}, which uses a stacked-hourglass architecture~\cite{newell2016stacked} and equivariance objectives to discover consistent keypoints. We simplify that design by removing the generative network used in prior work and by coupling the landmark detector directly to a thin-plate spline registrar~\cite{bookstein1989principal,keller2020thinplate}, trained with a multi-scale structural similarity loss~\cite{wang2003multiscale,wang2004image}. The result is a compact, unsupervised model that converges on the small sample sizes typical of spatial omics, that handles elastic-plus-rigid deformation through the thin-plate spline, and that can be instantiated once per modality so that multimodal datasets can be aligned in a shared latent space. Our contributions are fourfold:
\begin{itemize}
  \item A landmark-detection-first design that drops the generative branch of existing unsupervised detectors and couples the detector to a thin-plate-spline registrar, yielding consistent and stable landmarks on image counts typical of spatial omics experiments.
  \item A 3D adaptation in which landmarks are constrained to behave as anchor points for z-stack alignment, benchmarked against eight competing registration methods drawn from~\cite{kartasalo2018comparative} and from CODA~\cite{kiemen2022coda} on a 260-section mouse prostate dataset.
  \item A multimodal extension that trains separate detectors per modality and aligns their representations in a latent space, enabling joint registration of H\&E, Visium, and mass spectrometry data.
  \item A compatibility layer that allows ELD's landmarks to be plugged into existing spatial tools such as Eggplant~\cite{bergenstrahle2021eggplant}, replacing manual annotation and matching or exceeding it in accuracy.
\end{itemize}

\section{Related work}

\paragraph{Spatial omics and common coordinate frameworks.} Spatial transcriptomics methods such as the Visium array~\cite{stahl2016visualization}, in situ sequencing~\cite{ke2013insitu}, and recent atlases of the mouse central nervous system~\cite{chen2023atlas} and human heart~\cite{asp2019spatiotemporal} all require that sections be mapped to a shared anatomical frame before gene expression, cell type, or metabolite signals can be compared across animals or conditions. Museum-scale reviews~\cite{moses2022museum,marx2021method} highlight registration and coordinate frameworks as a recurrent bottleneck. Tools such as Eggplant~\cite{bergenstrahle2021eggplant} push the state of the art for landmark-based Visium alignment but depend on manual annotations, and STalign~\cite{clifton2023stalign} provides a diffeomorphic alternative that does not require landmarks but is less modular for multimodal settings.

\paragraph{Three-dimensional tissue reconstruction.} CODA~\cite{kiemen2022coda} reconstructs multi-centimetre tissues at cellular resolution by combining serial-section registration with deep-learning-based tissue labelling. The comparative benchmark of Kartasalo et al.~\cite{kartasalo2018comparative} evaluates seven reconstruction algorithms on a 260-section mouse prostate dataset with manually annotated landmarks, establishing the quantitative yardsticks -- Target Registration Error and Accumulated Target Registration Error -- that we adopt here. Classical non-rigid registration methods such as free-form deformation with B-splines~\cite{rueckert1999nonrigid} and thin-plate splines~\cite{bookstein1989principal} remain the workhorse of these pipelines.

\paragraph{Unsupervised landmark discovery in computer vision.} The unsupervised learning of object landmarks has been developed through equivariant labelling~\cite{thewlis2017unsupervised}, conditional image generation~\cite{jakab2018conditional}, structural representation learning~\cite{zhang2018unsupervised}, domain adaptation from a pretrained detector~\cite{sanchez2019object}, and self-supervised keypoint learning on unlabelled videos~\cite{jakab2020selfsupervised}. These methods rely heavily on the stacked hourglass backbone~\cite{newell2016stacked} and are typically evaluated on face benchmarks such as CelebA~\cite{liu2015celeba}, MAFL~\cite{zhang2016learning}, and AFLW~\cite{kostinger2011aflw}. While effective on images of identical objects, they assume dataset sizes and affine deformation regimes that do not hold in spatial biology.

\paragraph{Regularisation and augmentation for small data.} When training deep networks with only a handful of samples, regularisation is indispensable. Dropout~\cite{dropout2014hinton} has become the default tool for preventing co-adaptation of features, while classical augmentation techniques such as elastic deformation~\cite{simard2003best} remain surprisingly competitive on small image datasets. Our landmark dropout strategy and elastic augmentation pipeline directly build on these ideas. Image-quality losses such as SSIM~\cite{wang2004image} and its multi-scale extension MS-SSIM~\cite{wang2003multiscale} are more robust to the batch effects that plague histology, and we use MS-SSIM as the similarity term in all of our experiments.

\paragraph{Spatial omics analysis tooling.} We rely on the Scanpy framework~\cite{wolf2018scanpy} and its AnnData container~\cite{virshup2024anndata} for preprocessing and the Leiden algorithm~\cite{traag2019louvain} for clustering gene-expression neighbourhood graphs, following conventions established by Eggplant~\cite{bergenstrahle2021eggplant}.
